\section{Boxes with prescribed axes}\label{sec:boxes}

The mixed-boundary extension result lets us combine two holomorphic maps
on full disks, not just their germs. We then establish the parameter
dependence needed to combine maps from higher-dimensional polydisks.

\begin{proposition}[Disk axes]\label{prop:disk-axes}
Let \(\alpha,\beta:\D\to X\) be holomorphic maps with
\(\dd\alpha(T_\D)\subset E_1\), \(\dd\beta(T_\D)\subset E_2\), and
\(\alpha(0)=\beta(0)\). There is a unique holomorphic box
\(F:\D^2\to X\) such that
\[
 F(s,0)=\alpha(s),\qquad F(0,t)=\beta(t).
\]
\end{proposition}

\begin{proof}
Product coordinates at the common initial point give a unique box germ:
take the first coordinate from \(\alpha\) and the second from \(\beta\).
Any two boxes extending this germ agree on a connected common domain by
the identity theorem.

Fix an integer \(m\geq1\), and consider the convex domains
\[
 B_m(R)=\{(s,t):|s|^{2m}+|t|^{2m}<R^{2m}\},\qquad 0<R\leq1.
\]
Let \(R_*\) be the supremum of the radii to which the germ extends as
a box. The nested extensions agree and give a box on \(B_m(R_*)\).
Its values on the two axes are the prescribed maps, by the one-variable
identity theorem. Suppose that \(R_*<1\).

At a boundary point with \(st\neq0\), the complex tangent line has both
components nonzero: both coefficients of
\(\partial(|s|^{2m}+|t|^{2m})\) are nonzero there. Proposition~\ref{prop:mixed}
therefore extends the box across that point.
At a boundary point \((s_*,0)\), choose a disk \(U\) around \(s_*\),
with \(\overline U\subset\D\), such that \(\alpha(U)\) lies in a
product chart \(z=(z_1,z_2)\). Choose \(s'\in U\) with
\(|s'|<R_*\). For sufficiently small \(\epsilon>0\), the formula
\[
 G(s,t)=z^{-1}\bigl(z_1(\alpha(s)),z_2(F(s',t))\bigr),
 \qquad (s,t)\in U\times\D_\epsilon,
\]
is defined and is a box. Here \(\D_\epsilon=\{|t|<\epsilon\}\).
The product range of the chart makes the formula well-defined. Local
separation shows that \(G=F\) near \((s',0)\). Their common domain
\((U\times\D_\epsilon)\cap B_m(R_*)\) is convex, so they agree
throughout it. The other axis is treated symmetrically.

These boundary extensions glue. Indeed, restrict each one to an ordinary
Euclidean ball centered at its boundary point. If two such balls intersect,
their intersection contains a point of the segment joining their centers.
That segment lies in \(\overline{B_m(R_*)}\) by convexity. Since the
intersection is open, it consequently meets \(B_m(R_*)\). The two
extensions agree there, hence on their connected intersection. A finite
subcollection covers the compact boundary, giving an extension to a
neighborhood of \(\overline{B_m(R_*)}\). This neighborhood contains
\(B_m(R')\) for some \(R_*<R'<1\), a contradiction.

Thus the germ extends to \(B_m(1)\) for every \(m\). These domains
increase to \(\D^2\), and the resulting boxes agree on their connected
overlaps. Their union is the required box. The initial uniqueness
observation proves uniqueness on \(\D^2\).
\end{proof}

We will also use boxes with real parameters. A \emph{continuous box} on
a product of closed intervals is a continuous map into \(X\) that is
locally separated in product coordinates. Equivalently, every horizontal
slice locally follows an \(E_1\)-plaque and every vertical slice locally
follows an \(E_2\)-plaque. If a whole subrectangle maps into one product
chart, its first coordinate depends only on the horizontal parameter
and its second only on the vertical parameter: the unwanted coordinate
is locally constant along each connected slice.

\begin{lemma}[Uniqueness of continuous boxes]\label{lem:box-unique}
Two continuous boxes on \([0,1]^2\) with the same bottom and left edges
are equal.
\end{lemma}

\begin{proof}
Choose a Lebesgue number \(\lambda>0\) for a finite product-chart
cover of the compact metric space \(X\). Uniform continuity gives one
finite rectangular grid such that the image of each closed tile under
either map has diameter less than \(\lambda/3\). Whenever the two maps
agree at a tile's lower-left corner, the union of their tile images has
diameter less than \(2\lambda/3\), so both images lie in one product
chart. In that chart each box is determined by its bottom first
coordinate and left second coordinate. Starting with the prescribed
edges and proceeding by increasing sum of the tile indices proves
equality on the whole grid.
\end{proof}

\begin{proposition}[Polydisk axes]\label{prop:polydisk-axes}
Let \(S\) and \(T\) be unit polydisks of positive dimensions. Suppose
\(\alpha:S\to X\) and \(\beta:T\to X\) are holomorphic maps tangent
to \(E_1\) and \(E_2\), respectively, and
\(\alpha(p_0)=\beta(q_0)\) for some \((p_0,q_0)\in S\times T\).
There is a unique holomorphic box \(F:S\times T\to X\) with
\[
 F(p,q_0)=\alpha(p),\qquad F(p_0,q)=\beta(q).
\]
\end{proposition}

\begin{proof}
Polydisk automorphisms reduce the problem to \(p_0=q_0=0\).
For fixed \((p,q)\), the maps
\[
 \sigma\longmapsto\alpha(\sigma p),\qquad
 \tau\longmapsto\beta(\tau q)
\]
are holomorphic on disks of radii greater than one. After rescaling,
Proposition~\ref{prop:disk-axes} gives a box whose restriction to the
closed real square is denoted by \(R_{p,q}\). Set
\[
 F(p,q)=R_{p,q}(1,1).
\]
The constant-axis cases and Lemma~\ref{lem:box-unique} give the prescribed
axis values. It remains to prove holomorphicity and the two tangencies;
pointwise existence of \(R_{p,q}\) alone does not give parameter
continuity.

Fix \((p_*,q_*)\), and choose a finite grid whose closed tile images
under \(R_{p_*,q_*}\) lie in product charts. We reconstruct this square
for nearby \((p,q)\), using these fixed charts. Begin with its two
given axes, and process tiles from the bottom and left. The induction
maintains the following properties: each constructed tile depends
jointly continuously on its real parameters and on \((p,q)\); it is a
continuous box in the real parameters; and, at each fixed pair of real
parameters, it is holomorphic in \((p,q)\), with the \(p\)-directions
in \(E_1\) and the \(q\)-directions in \(E_2\). These properties hold
on the initial axes.

To fill one tile, denote its already constructed bottom and left edges
by \(B(\sigma;p,q)\) and \(L(\tau;p,q)\). At the reference parameter
their images lie in the selected product chart \(z=(z_1,z_2)\).
Joint continuity and compactness of those edges put both images in that
same chart after shrinking the parameter neighborhood. Define the tile
by
\begin{equation}\label{eq:parameter-tile}
 z^{-1}\bigl(z_1(B(\sigma;p,q)),z_2(L(\tau;p,q))\bigr).
\end{equation}
The edges meet at a common corner. Plaque constancy on each edge shows
that this formula reproduces both of them. It is jointly continuous and
is separated in \(\sigma,\tau\). For each fixed \(\sigma,\tau\), it
is holomorphic in \((p,q)\); moreover, the induction hypothesis gives
\[
 \partial_q\bigl(z_1\circ B\bigr)=0,
 \qquad
 \partial_p\bigl(z_2\circ L\bigr)=0.
\]
Thus the two parameter tangencies are preserved as well. At
\((p_*,q_*)\), the reconstructed tile is the original tile, so the
induction continues. There are only finitely many tiles and therefore
only finitely many shrinkings.

The tile maps agree on common edges and form a continuous box across
the seams. To check the latter assertion at a seam or a vertex, use a
product chart around its common image. On a sufficiently small real
neighborhood, the unwanted coordinate along a horizontal or vertical
slice is constant on each tile piece and agrees at their endpoints;
it is therefore constant across the seam. By
Lemma~\ref{lem:box-unique}, the reconstructed square is \(R_{p,q}\)
for every nearby parameter pair. Its upper-right value is consequently
\(F(p,q)\), proving that \(F\) is a holomorphic box near
\((p_*,q_*)\), and hence everywhere. Finally, the axes determine its
germ at \((0,0)\), so the identity theorem proves uniqueness.
\end{proof}
